\section{Discussion: policy interpretation and limits}\label{sec:institutions}

The policy lesson concerns the counterfactual used to appraise free basic access, not an unconditional case for public entry. Where a paid specialist screens founders, the relevant comparison is welfare under private-only provision against welfare after free basic access is added alongside a paid premium. Additional access restores direct service value but can displace premium use and induce ventures with negative local commercialisation incidence. Proposition~1 identifies when the fee-induced access effect is strong enough that greater skill complementarity raises the conditional public cost ceiling despite such adverse incidence. Governments must still compare that ceiling with the real fixed cost; a national evaluation additionally includes the partner region's welfare effect, which can have either sign.

This distinction suggests a policy diagnostic: identify eligible firms and those excluded at the private-only fee; separate basic-service benefits from premium use; and locate consumer, incumbent and founder gains or losses from venture entry. Local fee payments are transfers, whereas service operating costs use real resources. Registrations, matched-project counts and observed success rates do not identify these welfare components or the causal private-fee response. The measurement and identification requirements are discussed in Section~\ref{sec:empirical}.

Institutional examples motivate these questions without establishing the proposition empirically. The European Digital Innovation Hubs network supports access to experimentation and expertise, with free or discounted services valued against market prices \citep{EuropeanCommissionEDIH,EDIHFAQ2026}. In the British Catapult debate, parliamentary concerns about competing for projects and funding contrast with the government's defence of mixed financing and commercial revenues \citep{UKCommittee2026,UKGovernment2026}; neither position provides causal evidence of crowding in or out, and Catapults do not map literally to separate public and private providers. The US Manufacturing Extension Partnership's reported help connecting a Montana inventor to manufacturing capabilities illustrates invention--implementation matching, not a verified fee-screening response \citep{NISTFiberReaper2024}.

The restrictions are consequential. Proposition~1 assumes uniform founder ability perfectly aligned with premium valuation, an interior local fee, full partner-region participation and sufficiently weak matching. Premium advice does not directly raise technical success, and founders are randomly paired rather than optimally matched. The finite-matching certificate covers a fixed continuous parameter region, not arbitrary environments. The four-party configuration switch is confined to a narrow public-cost interval and disappears under the reported Beta$(2,2)$ alternative; the recorded adverse Bertrand test is not a newly established general theorem. No general four-player comparative static or quantitatively important effect in an observed region follows. Finally, the government supplies taxpayer-financed basic intermediation, not necessarily a building: its location, capital costs, capacity, taxation, subsidies, procurement and remote alternatives lie outside the policy menu. The paper identifies an appraisal boundary, not the optimal instrument or a construction recommendation.
